\section{Code and Model}
We open-source our code and model as follows:

\faGithub\ \textbf{\textcolor{blue}{Code:} \href{https://github.com/LCM-Lab/Elastic-Attention}{https://github.com/LCM-Lab/Elastic-Attention}}

\faDatabase\ \textbf{\textcolor{blue}{Model:} \href{https://modelscope.cn/collections/LCM\_group/Elastic-Attention}{https://modelscope.cn/collections/LCM\_group/Elastic-Attention}}

\section{Related Work}
\label{sec:related_work}

\subsection{Sparse Attention Mechanisms}
\label{subsec:train_ssa}
To mitigate the quadratic complexity of standard attention, existing research has broadly evolved along two trajectories: inference-time heuristics and training-aware sparsification.
Inference-time heuristics typically employ static patterns, such as fixed sliding windows or strides~\citep{streamingllm, TriangleMix, Beltagy2020Longformer}, to restrict the receptive field. 
To capture more dynamic dependencies, content-aware approaches have been proposed. Token eviction policies discard uninformative tokens based on accumulated importance scores~\citep{h2o, snapkv, scissorhands}, while kernel-based estimators identify salient blocks to bypass redundant computations~\citep{jiang2024minference}. 
Complementarily, prefill optimizers leverage importance-driven selection to accelerate long-context processing~\citep{flexprefill, xattention, spargeattn, peng2025acceleratingprefillinglongcontextllms, ji2025salelowbitestimation}. Despite their effectiveness, these heuristic methods often hinge on sensitive hyperparameters, limiting their robustness across varying tasks.

In contrast, internalize sparsity within the optimization objective to align training with sparse inference.
A primary direction involves learnable selection. For instance, SeerAttention~\citep{gao2024seerattention}, NSA~\citep{NSA} and MoBA~\citep{moba} employ learnable gates and hierarchical constraints to approximate ground-truth attention patterns. 
To bridge the gap between dense pre-training and sparse adaptation, InfLLM-v2~\citep{infllmv2} introduces a dense-sparse switchable mechanism via parameter-free pooling, while DSA~\citep{deepseekai2024deepseekv32} utilizes a lightning indexer with a two-stage training strategy to efficiently filter top-$k$ key-value pairs.
However, most of these methods focus on fine-grained block-level selection within a fixed attention mechanism, rather than dynamically adapting the attention mode itself based on input complexity.

\subsection{Hybrid Efficient Architectures}
\label{subsec:hybrid_arch}
To balance efficiency with performance, hybrid architectures strategically integrate Full Attention (FA) with linear-complexity operators.
The dominant paradigm, inter-layer hybridization, interleaves linear layers (e.g., SSMs or RNNs) with standard attention layers to recover associative recall capabilities~\citep{ku2025systemsalgorithmsconvolutionalmultihybrid,dao2024transformersssmsgeneralizedmodels}. Notable large-scale implementations, such as Jamba~\citep{lieber2024jambahybridtransformermambalanguage}, utilize fixed block-wise ratios, while variants optimize memory via shared global blocks~\citep{glorioso2024zambacompact7bssm} or sliding windows~\citep{ren2024samba}.

More recently, intra-layer hybridization has emerged to refine granularity. For example, PruLong~\citep{bhaskar2025cache} and DuoAttention~\citep{duoattention} combine FA and Streaming Sparse Attention (SSA) within individual layers, assigning different heads to different modes. LongCat~\citep{zhang2025efficient} proposes the LoZA mechanism, constructing a static ZigZag topology by replacing low-sensitivity Multi-head Latent Attention (MLA) modules with linear-complexity SSA.
Nevertheless, a critical limitation of these approaches is their reliance on static topologies or pre-defined ratios determined prior to inference. Such rigid designs lack the flexibility to distinguish between sparsity-robust and sparsity-sensitive tasks dynamically, often leading to suboptimal resource allocation for diverse inputs.